\prefacesection{Abstract}

Autonomous underwater vehicles (AUVs) are constrained by low-bandwidth communication channels, as well as limited onboard computation. Current approaches either rely heavily on cloud processing or transmitting raw images, which are infeasible on untethered AUVs. In this work, we introduce ATHENA: AI Transformer for Hydro-Environmental Navigation and Analysis, an edge-capable Vision-Language Model (VLM) pipeline that can enable AUVs to describe underwater scenes as text, allowing for efficient data transmission and real-time situational awareness. To address common underwater challenges such as color distortion, light attenuation, and haze, we apply a multi-step preprocessing pipeline that enhances image clarity for the VLM. By leveraging 4-bit quantization, ATHENA reduces the memory footprint to just 652.86 MB, making it deployable on embedded hardware. Preliminary results show concise, accurate captions with significant improvements on standard image captioning benchmarks such as BLEU, ROUGE-L, and CIDEr.

% Tips on how to write abstract:
% https://communities.springernature.com/posts/how-to-write-an-abstract
% https://unl.libguides.com/c.php?g=51569&p=2633458